\section{The starting cases}\label{ext:starting-cases}
\subsection{The first three sample sizes}
\label{ext:small-foundations}

The small cases provide quantitative starting bounds for the later
induction.  The second and third cases use the same idea: a hypothetical
violating extremizer has bounded support, and a polynomial majorizes its
closed sum event on the whole support product.  Independence then evaluates
the polynomial using moments of the same marginal.  No restriction on the
number of atoms is imposed.

\begin{theorem}\label{ext:n123}
For every standardized law with finite third absolute moment,
\[
 C_1\le\frac25,\qquad C_2\le\frac{2\sqrt2}{7},\qquad
 C_3\le\frac{59\sqrt3}{250}.
\]
All three upper bounds are strictly smaller than $\CE$.
\end{theorem}

\paragraph{One summand.}
The certificate consists of 150 consecutive rational cells $[a,b]$
covering $[-2,1]$, with two rational coefficients $u,v$ in each cell.
It proves
\begin{equation}\label{ext:n1-dual}
 L_a-v+uy+vy^2+\frac25|y|^3-\ind_{\{y\le b\}}\ge0
 \quad(y\in\R),\qquad L_a\le\Phi(a).
\end{equation}
Taking expectations cancels the linear and quadratic moment terms.
For $t\in[a,b]$ this gives
\[
 F_P(t)-\Phi(t)\le F_P(b)-L_a\le\frac25\rho.
\]
For $t\le-2$, Cantelli gives $F_P(t)\le1/(1+t^2)\le1/5$.
For $t\ge1$, the discrepancy is at most $1-\Phi(1)<2/5$.
Since $\rho\ge1$, these bounds cover both exterior rays.
Reflection gives the negative branch, including atoms.

For completeness, the verification of \eqref{ext:n1-dual} splits the
real line at $-10,0,b,10$.  On each finite interval the expression is a
cubic, checked in the Bernstein basis with subdivision; on either
infinite ray its translated coefficients are nonnegative.  The value at
$b$ uses the included, closed-event polynomial.  The normal lower bound
uses Machin's identity, integer square-root brackets and an alternating
Gaussian series.  All 150 cells and 1,128 terminal Bernstein intervals
pass exact rational verification.

\paragraph{A common support argument for two and three summands.}
Let $K_2=2\sqrt2/7$, $K_3=59\sqrt3/250$ and suppose $C_n>K_n$,
where $n=2$ or $3$.  Corollary~\ref{ext:ratio-attainment} supplies a
positive closed-CDF maximizing pair.  Write
\[
 z=H/\sqrt n,\quad F=\Pp(S_n\le H),\quad D_0=F-\Phi(z),
 \quad \lambda=D_0/\rho,\quad \beta=\E[X|X|].
\]
Here $D_0$ denotes the discrepancy, not the support-gap defect used in
the main argument.  The contact condition is equivalently
\begin{align}
 n\Pp(S_{n-1}\le H-y)&\le Q_n(y),\notag\\
 Q_n(y)&=\lambda|y|^3-\frac{z\phi(z)+3D_0}{2}y^2
 -(\sqrt n\phi(z)+3\lambda\beta)y
 +nF+\frac{z\phi(z)+D_0}{2},\label{ext:n23-contact}
\end{align}
with equality on the entire closed support.
The published affine estimate
\begin{equation}\label{ext:n23-affine}
 \frac{\sqrt n\,\Delta_n(P)}{\rho(P)}\le\frac{1661}{5000}
       \left(1+\frac{429}{1000\rho(P)}\right)
\end{equation}
from \citet{Shevtsova2013} rules out $\rho\ge2$ at $n=2$ and
$\rho\ge15/8$ at $n=3$.  Together with $|\beta|\le1$,
$\phi(z)<2/5$ and $|z\phi(z)|<1/4$, it gives the following lower
bounds $Q_n(y)\ge P_n(|y|)$:
\[
\begin{array}{c|c|c}
 n&\rho\text{ ceiling}&P_n(r)\\\hline
 2&2&\displaystyle\frac27r^3-\frac{67}{56}r^2-\frac{23}{14}r-\frac18\\[2pt]
 3&15/8&\displaystyle\frac8{35}r^3-\frac45r^2-\frac85r-\frac18
\end{array}
\]
Indeed, for $n=2$, $\lambda>2/7$, $D_0<5/7$ and
$\lambda<5/14$ suffice: $\sqrt2<10/7$ gives
$\sqrt2\phi(z)+3\lambda<4/7+15/14=23/14$.
For $n=3$, the corresponding bounds are
$\lambda>8/35$, $D_0<9/20$ and $\lambda<2/7$.
In the latter case \eqref{ext:n23-affine} gives $C_3\rho<77/100$,
and $\sqrt3>12/7$ then gives $D_0=C_3\rho/\sqrt3<9/20$.
The shifted differences are
\begin{align*}
 P_2(6+s)-2&=\frac27s^3+\frac{221}{56}s^2+
                  \frac{104}{7}s+\frac{373}{56},\\
 P_3(6+s)-3&=\frac8{35}s^3+\frac{116}{35}s^2+
                  \frac{472}{35}s+\frac{2197}{280}.
\end{align*}
Every coefficient is positive.  Equality in
\eqref{ext:n23-contact} therefore excludes every support point with
$|y|\ge6$, because its left side is at most $n$.

\paragraph{The two-summand majorants.}
Put $\gamma=\E X^3$ and
\[
 v=(1,y,y^2,y^3,|y|^3),\qquad \mu=\E v(X)=(1,0,1,\gamma,\rho).
\]
For a rational symmetric coefficient array indexed by $i\le j$, let
\[
 M_c(x,y)=\sum_i c_{ii}v_i(x)v_i(y)+
 \sum_{i<j}c_{ij}\{v_i(x)v_j(y)+v_j(x)v_i(y)\},
\]
and $q_c(\gamma,\rho)=\sum_i c_{ii}\mu_i^2+
2\sum_{i<j}c_{ij}\mu_i\mu_j$.
Independence gives $\E M_c(X_1,X_2)=q_c(\gamma,\rho)$.
Theorem~\ref{ms:sharp-third} puts the actual moments in
\[
 \mathcal D_2=\{(\gamma,\rho)\in[-2,2]\times[1,2]:
 \rho^2(\gamma^2+4)-(\gamma^2+2)^2\ge0\}.
\]
The 70 certificate rows use
$a_j=-5/2+j/20$, $b_j=a_j+1/20$, $0\le j<70$.
Their exact inequalities are
\begin{align*}
 M_{c^{(j)}}(x,y)&\ge\ind_{\{x+y\le b_j\}}
                    &&((x,y)\in[-6,6]^2),\\
 q_{c^{(j)}}(\gamma,\rho)&\le\Phi(a_j/\sqrt2)+\frac27\rho
                    &&((\gamma,\rho)\in\mathcal D_2).
\end{align*}
Each sign quadrant is cut by the event line and triangulated into
closed triangles.  Degree-six Bernstein coefficients, with exact
subdivision, establish the first inequality.  A second checker
independently homogenizes the original polynomial on every triangle
and checks integer coefficients.  Both checkers accept all 70 rows:
35,878 triangle nodes and 134 moment boxes.  The second inequality is
checked on moment rectangles; a rectangle is discarded only if
$\rho_+^2(g_*+4)-(g_*+2)^2<0$, where $g_*$ is the minimum of
$\gamma^2$ on that rectangle.  This is valid because the expression
increases with $\rho$ and decreases with $\gamma^2$ when $\rho\le2$.

For a threshold in $[a_j,b_j]$, expectation of the first inequality and
the second give discrepancy at most $2\rho/7$.
For $H\le-5/2$, Cantelli gives $2/(2+H^2)\le8/33<2/7$;
for $H\ge1$, $1-\Phi(H/\sqrt2)<2/7$.
These bounds contradict the chosen violating maximizer.  This proves
the second assertion of Theorem~\ref{ext:n123}.

\paragraph{The three-summand majorants.}
Use the features
\[
 (v_0,\ldots,v_6)(y)=(1,y,y^2-1,y^3,|y|^3,y^4,y^5).
\]
Let $\mathcal I$ be the 59 nondecreasing triples from $\{0,\ldots,6\}$
whose largest entry is less than 5 or which contain an entry 1 or 2.
For $\alpha\in\mathcal I$, average the associated product over its
distinct coordinate permutations, obtaining $V_\alpha$.
Set $M_c=\sum_\alpha c_\alpha V_\alpha$.
The bounded-support argument above makes every feature integrable.
Every term involving a fourth or fifth power contains a zero-mean
factor $X_i$ or $X_i^2-1$ in a different coordinate.  Independence
therefore gives the exact identity
\[
 \E M_c(X_1,X_2,X_3)=q_c(\gamma,\rho)
 =\sum_{\substack{\alpha\in\mathcal I\\\alpha\text{ uses only }0,3,4}}
     c_\alpha\gamma^{\#3(\alpha)}\rho^{\#4(\alpha)}.
\]
There are ten surviving monomials and 49 zero-expectation controls.
The controls improve the pointwise majorant without adding moment
parameters.

The 111 rational rows form a consecutive closed cover of
$[-7/2,3/2]$.  They satisfy
\begin{align*}
 M_{c^{(j)}}(x_1,x_2,x_3)&\ge\ind_{\{x_1+x_2+x_3\le b_j\}}
                    &&((x_1,x_2,x_3)\in[-7,7]^3),\\
 q_{c^{(j)}}(\gamma,\rho)&\le\Phi(a_j/\sqrt3)+\frac{59}{250}\rho
                    &&((\gamma,\rho)\in\mathcal D_3),
\end{align*}
where $\mathcal D_3$ replaces $2$ by $15/8$ in the rectangular
bounds defining $\mathcal D_2$.
Symmetry reduces the sign patterns to four ordered-coordinate
orthants.  A box is discarded only when it has no ordered point.
On every other box, degree-five tensor Bernstein bounds establish
the pointwise inequalities.  On a box crossing the event plane,
the included side is checked by $M-1\ge s(x_1+x_2+x_3-b_j)$ with
$s\le0$, and the other side by the analogous inequality with $s\ge0$.
Existence of a slope is checked by intersecting exact rational intervals;
zero plane coefficients are checked directly.  This retains the closed
event boundary.

Fresh exact verification accepts all 111 cells, with 501,156 cube
nodes, 219,860 terminal leaves and 349 moment boxes.  An additional
95,904 exact interpolation identities compare the initial polynomial
expansion with direct evaluation of the original features.  Those
identities check the expansion; the cube verification proves positivity.
As in the two-summand case, expectation bounds the discrepancy on each
threshold interval.  The exterior bounds are
$3/(3+H^2)\le12/61<59/250$ for $H\le-7/2$ and
$1-\Phi(H/\sqrt3)<59/250$ for $H\ge3/2$.
They contradict the violating maximizer and prove the third assertion.
The first assertion was proved above.  Rational square-root and Machin
bounds give $K_2<0.405<0.409<\CE$ and $K_3<0.409<\CE$.

\subsection{The fourth case: localization before numerical comparison}
\label{ext:n4-localization-section}

The third bound is needed quantitatively.  Merely knowing $C_3\le\CE$
would not give the following support reduction.
Write
\[
 K_3=\frac{59\sqrt3}{250},\quad c_0=\frac{256}{625},\quad
 T=\frac{4097}{10000};\qquad K_3<c_0<T<\CE.
\]
For any standardized law set
$\beta=\E[X|X|]$, $u=\E X_+=\E X_->0$,
$v_\pm=(1\pm\beta)/2$, and
\begin{equation}\label{ext:n4-reference}
 a=\frac{v_-}{u},\quad b=\frac{v_+}{u},\quad
 w=\frac{4u^2}{1-\beta^2},\quad
 Q=\frac ua\delta_{-a}+(1-w)\delta_0+\frac ub\delta_b.
\end{equation}
Cauchy's inequality gives $0<w\le1$ and makes $Q$ a standardized
probability law.  Put $r=\rho(Q)=(1+\beta^2)/\sqrt{w(1-\beta^2)}$.
\begin{lemma}\label{ext:n4-excess}
The nonnegative quantity
\[
 W(X)=X_+(X_+-b)^2+X_-(X_--a)^2
\]
has expectation $\eta=\rho(P)-r$.  Thus $r\le\rho(P)$, with
equality exactly when $P=Q$.  Moreover, writing $c(X)=b$ for $X>0$
and $c(X)=a$ for $X<0$,
\begin{align}
 \E X^3&=b-a+\E[\operatorname{sgn}(X)W(X)],\notag\\
 \E X^4&=a^2-ab+b^2+\E[(|X|+2c(X))W(X)]
 \label{ext:n4-physical-moments}
\end{align}
whenever the fourth moment is finite.
\end{lemma}
\begin{proof}
On either half-line put $y=|X|$ and $c=a$ or $b$.  The identities
$y^3=y(y-c)^2+2cy^2-c^2y$ and
$y^4=(y+2c)y(y-c)^2+3c^2y^2-2c^3y$ integrate using
$\E[y;\text{half-line}]=u$ and
$\E[y^2;\text{half-line}]=uc$.
If $\E W=0$, the support is contained in $\{-a,0,b\}$ and the
first two moments fix the masses.  Conversely $P=Q$ has $W=0$.
\end{proof}
The three-point law is a reference;
the original law remains arbitrary.

\begin{proposition}\label{ext:n4-localization}
If $C_4>c_0$, a positive closed-CDF maximizing pair satisfies
\[
 \supp P\subset(-101/25,101/25),\quad\rho(P)<46/25,
 \quad 625/2116<w\le1,\quad|\beta|<321/500.
\]
In particular it belongs to the closed numerical class with
$|X|\le17/4$ and $|\beta|\le13/20$.
\end{proposition}
\begin{proof}
Use Corollary~\ref{ext:ratio-attainment}.  The affine inequality
\eqref{v9:entry:structural} first gives
$\rho<979013/532500<46/25$, without a support assumption.
Since $r\le\rho$ and $w\le1$, the inequality
\[
 \frac{(1+(321/500)^2)^2}{1-(321/500)^2}>(46/25)^2
\]
gives $|\beta|<321/500$.
Put $C=C_4$, $z=H/2$, $s=\sqrt{4/3}$.  The contact expression,
using the quantitative bound on the three-summand predecessor, obeys
\begin{align*}
 I_{\rm new}(y)\le{}&\mathcal H_4(z,y)+4\rho(K_3-C)
                    +4\rho(s-1)K_3\\
 &+C\{-|y|^3+\tfrac32\rho y^2+3\beta y-\rho/2\},
\mathcal H_4(z,y)={}&8\!\left\{\Phi\left(\frac{z-y/2}{\sqrt{3/4}}\right)-\Phi(z)\right\}
                 \\
 &+4y\phi(z)+z\phi(z)(y^2-1).
\end{align*}
Taylor expansion in the translation and the variance gives, uniformly
in $z,y$,
\[
 |\mathcal H_4(z,y)|<\frac{\phi(0)}6|y|^3+\frac{|y|}{5}+\frac1{15}.
\]
Here $\sup|\phi''|=\sup|(x\phi(x))'|=\phi(0)<2/5$ and
$\sup|x(3-x^2)\phi(x)|<3/5$ suffice.  For the last bound, on
$t=x^2\le3$ compare $e^t$ with its cubic Taylor polynomial and use
\[
 \tfrac94(1+t+t^2/2+t^3/6)-t(3-t)^2
 =\tfrac34\{7(t-9/14)^2+3/28\}+\tfrac58t^2(3-t)>0.
\]
For $t\ge3$, use $|x|^3\phi(x)<1/2$.
The second derivative of
$\Phi((z-y/2)/\sqrt{1-\lambda})$ has absolute value below $4/15$
on $0\le\lambda\le1/4$, giving the remainder $1/15$ after
multiplication by eight.

Discard $4\rho(K_3-C)<0$ and use $4(s-1)<5/8$.
For $R=|y|\ge4$, the remaining coefficient of $C$ is at most
$B(R)=-R^3+69R^2/25+963R/500+23/100$.
All coefficients of $B(4+t)$ are negative, so substituting $C>c_0$
is now valid.  Consequently
\[
 I_{\rm new}(y)<-\frac{643}{1875}R^3+
 \frac{17664}{15625}R^2+\frac{77257}{78125}R+\frac{7541}{46875}.
\]
The coefficients after translation by $R=101/25+t$ are all negative.
This excludes $|y|\ge101/25$ from the closed support by contact.
Finally $r^2\ge1/w$ gives $w>625/2116$.
\end{proof}

% The numerical ROOT4 physical-class theorem and its full common-measure,
% quadrature, low/high-excess and closed-cover proof must follow here.
% Do not infer C4<=T merely from Proposition ext:n4-localization.
